% TOP_PROBLEM: {"id": 2302045, "title": "Research Problems in Function Theory — Problem 2.45", "category_id": 8, "category": {"id": 8, "name": "analysis", "display_name": "Analysis", "description": "Limits, continuity, calculus, and function theory.", "slug": "analysis", "order_index": 8, "created_at": "2026-07-31T15:26:25.671Z"}, "status": "open", "problem_number": "AMR-022-2045", "set_id": 13}
% FIRST_POSTED: 2026-10-02
% ENTRY_KIND: statement_only
% UnsolvedMath ID 2302045; source catalogue code AMR-022-2045
% !TeX program = lualatex
% Definition and sources only; this file is not an LLM solution attempt.
\documentclass[11pt,a4paper]{article}
\usepackage[margin=25mm]{geometry}
\usepackage{fontspec}
\setmainfont{lmroman10-regular.otf}[BoldFont=lmroman10-bold.otf,
 ItalicFont=lmroman10-italic.otf,BoldItalicFont=lmroman10-bolditalic.otf]
\usepackage{amsmath,amssymb,mathtools}
\usepackage{unicode-math}
\setmathfont{latinmodern-math.otf}
\AtBeginDocument{\let\setminus\smallsetminus}
\usepackage{xurl,hyperref}
\hypersetup{colorlinks=true,urlcolor=blue}
\setcounter{secnumdepth}{0}
\setlength{\parindent}{0pt}
\setlength{\parskip}{5pt}
\setlength{\emergencystretch}{3em}
\begin{document}
\section*{Research Problems in Function Theory — Problem 2.45}\label{2302045}
{\small UnsolvedMath ID 2302045 · AMR-022-2045 · Analysis · AMR Open Problem Lists}
\subsection{Definitions and mathematical statement}
The Bessel function of the first kind and order zero is the entire even function
\[
J_0(z)=\sum_{k=0}^{\infty}\frac{(-1)^k}{(k!)^2}\left(\frac z2\right)^{2k}.
\]
For each angle $\theta\in[0,2\pi)$ consider the open ray
$R_\theta=\{te^{i\theta}:t>0\}$. Is it true that
\[
\#\{t>0:J_0(te^{i\theta})=1\}\le1
\qquad\text{for every }\theta?
\]
The count is of distinct nonzero solutions, not their analytic multiplicities. The common vertex is excluded because $J_0(0)=1$ automatically; including it would change the intended ray question. Equivalently, can two distinct nonzero roots of $J_0(z)-1$ ever have the same argument modulo $2\pi$? The evenness of $J_0$ pairs roots on opposite rays and does not violate the proposed assertion.
The question asks for a global statement covering all arguments and all moduli. Large-argument asymptotic estimates alone, even when they bound the number of roots on a fixed ray, do not decide uniqueness among the remaining finite roots. The source connects this assertion to removing an exceptional set in a theorem of Delsarte and Lions.
\subsection{Short English statement}
Can the equation $J_0(z)=1$ have two distinct nonzero solutions on the same ray from the origin?
\subsection{Sources}
\begin{itemize}
\item[S1] ULAM.AI and the UnsolvedMath Contributors, supplied problems.json, record 2302045 (AMR-022-2045), AMR Open Problem Lists. Catalogue identity and source transcription; CC BY 4.0.
\url{https://huggingface.co/datasets/ulamai/UnsolvedMath}.
\item[S2] Hayman - Research Problems in Function Theory (2018), Problem 2.45. Original source indexed by AMR-022-2045.
\url{https://arxiv.org/abs/1809.07200}.
\end{itemize}
\end{document}
